% !TEX root = ../main/main.tex
% Section: Semantic Stack for LLM-Powered Feature Engineering
% Created: December 21, 2025
% Purpose: Document LLM integration for semantic feature extraction

\section{Semantic Stack: LLM-Powered Feature Engineering}
\label{sec:semantic-stack}

\subsection{Motivation: Beyond Quantitative Features}

The 342-feature quantitative pipeline documented in Chapters~3--4 captures structured performance metrics: EPA, success rates, NextGen tracking data, and derived statistics. However, NFL betting markets also respond to \textit{unstructured} information: injury reports, social media sentiment, coaching commentary, and media narratives. This section documents the ``Semantic Stack''---a layered LLM architecture designed to extract actionable features from unstructured text sources.

\paragraph{The Information Gap.}
Consider a typical injury report: ``Patrick Mahomes questionable with ankle injury.'' A human bettor immediately processes:

\begin{itemize}[itemsep=2pt,parsep=0pt]
    \item \textbf{Entity recognition:} Mahomes = elite QB, ankle = mobility-affecting
    \item \textbf{Status interpretation:} ``Questionable'' = 50-75\% chance of playing, but if playing, likely limited
    \item \textbf{Market impact:} Line should move 2-4 points if ruled out; 0.5-1 point if limited
    \item \textbf{Historical pattern:} Mahomes has played through injuries before (2022 AFC Championship)
\end{itemize}

Traditional feature engineering cannot capture this reasoning chain. The Semantic Stack bridges this gap through structured LLM inference.

\paragraph{Expected Impact (Pre-Validation Targets).}
Based on prior work in sports analytics NLP~\cite{Ramirez2023} and market sentiment analysis~\cite{Chen2024}, we originally targeted:

\begin{itemize}[itemsep=2pt,parsep=0pt]
    \item \textbf{Win rate uplift:} +1.5\% to +2.5\% (from 51.0\% baseline toward 52.5--53.5\%)
    \item \textbf{Props MAE reduction:} -5 to -10 yards
    \item \textbf{Value bet identification:} 5-10 additional opportunities per week
\end{itemize}

\textbf{Critical update:} A/B testing on 2025 holdout data (Section~\ref{sec:semantic-ab-test}) showed semantic features \emph{underperformed} the baseline (58.5\% → 51.8\%, $p=0.21$). These targets remain aspirational pending Phase 2 improvements to NER status extraction and feature selection.

\subsection{Architecture Overview}

The Semantic Stack implements a four-layer architecture, from data ingestion through decision support (Figure~\ref{fig:semantic-stack-architecture}).

\begin{figure}[htbp]
\centering
\begin{tikzpicture}[
    node distance=1.0cm,
    layer/.style={rectangle, draw=accentblue, fill=accentblue!10, rounded corners, minimum height=1.8cm, minimum width=10cm, font=\footnotesize},
    component/.style={rectangle, draw=black, fill=white, rounded corners, minimum height=0.7cm, font=\scriptsize, inner sep=4pt}
]

% Layer 4: Decision Support
\node[layer] (layer4) {};
\node[font=\footnotesize\bfseries, accentblue] at ([yshift=0.5cm]layer4.center) {L4: Decision Support};
\node[component] at ([xshift=-2.8cm, yshift=-0.3cm]layer4.center) (vbd) {ValueBetDetector};
\node[component] at ([yshift=-0.3cm]layer4.center) (expl) {LLM Explanations};
\node[component] at ([xshift=2.8cm, yshift=-0.3cm]layer4.center) (hitl) {Human-in-Loop};

% Layer 3: Feature Engineering
\node[layer, below=1.3cm of layer4] (layer3) {};
\node[font=\footnotesize\bfseries, accentblue] at ([yshift=0.5cm]layer3.center) {L3: Semantic Features (26)};
\node[component] at ([xshift=-3cm, yshift=-0.3cm]layer3.center) (sent) {Sentiment};
\node[component] at ([xshift=-1cm, yshift=-0.3cm]layer3.center) (inj) {Injury Severity};
\node[component] at ([xshift=1cm, yshift=-0.3cm]layer3.center) (media) {Media Buzz};
\node[component] at ([xshift=3cm, yshift=-0.3cm]layer3.center) (match) {Matchup};

% Layer 2: LLM Inference
\node[layer, below=1.3cm of layer3] (layer2) {};
\node[font=\footnotesize\bfseries, accentblue] at ([yshift=0.5cm]layer2.center) {L2: LLM Processing};
\node[component] at ([xshift=-3cm, yshift=-0.3cm]layer2.center) (ollama) {Ollama};
\node[component] at ([xshift=-1cm, yshift=-0.3cm]layer2.center) (mlx) {MLX};
\node[component] at ([xshift=1cm, yshift=-0.3cm]layer2.center) (gemini) {Gemini};
\node[component] at ([xshift=3cm, yshift=-0.3cm]layer2.center) (claude) {Claude};

% Layer 1: Data Ingestion
\node[layer, below=1.3cm of layer2] (layer1) {};
\node[font=\footnotesize\bfseries, accentblue] at ([yshift=0.5cm]layer1.center) {L1: Data Ingestion};
\node[component] at ([xshift=-3cm, yshift=-0.3cm]layer1.center) (espn) {ESPN News};
\node[component] at ([xshift=-1cm, yshift=-0.3cm]layer1.center) (twitter) {Twitter/X};
\node[component] at ([xshift=1cm, yshift=-0.3cm]layer1.center) (injury) {Injuries};
\node[component] at ([xshift=3cm, yshift=-0.3cm]layer1.center) (db) {84K+ Records};

% Arrows
\draw[->, thick, accentblue] (layer1) -- (layer2);
\draw[->, thick, accentblue] (layer2) -- (layer3);
\draw[->, thick, accentblue] (layer3) -- (layer4);

\end{tikzpicture}
\caption{Semantic Stack Architecture. Four-layer design from data ingestion (L1) through decision support (L4). Layer 2 implements fallback chain for inference reliability: local models (Ollama, MLX) preferred, API fallback (Gemini, Claude) for complex queries or availability issues. Layer 3 generates 38 semantic features (17 per team $\times$ 2 + 4 matchup-level) integrated into existing GameFeatureExtractor pipeline. Layer 4 provides human-readable explanations for value bet recommendations.}
\label{fig:semantic-stack-architecture}
\end{figure}

\subsection{Layer 1: Data Ingestion Pipeline}

\paragraph{News Sources.}
The \texttt{NewsIngester} module (\texttt{nfl\_analytics/semantic/\allowbreak ingestion/\allowbreak news\_ingester.py}) aggregates content from:

\begin{itemize}[itemsep=2pt,parsep=0pt]
    \item \textbf{ESPN API:} Official news, injury reports, game previews (primary source)
    \item \textbf{Team RSS feeds:} Official team announcements
    \item \textbf{Beat reporter Twitter:} Verified NFL insiders (Schefter, Rapoport, etc.)
\end{itemize}

Each article is classified by type (injury, trade, game preview, general) using keyword matching with word-boundary enforcement to prevent false positives (e.g., ``linebacker'' falsely matching ``back'' injury keyword).

\paragraph{Social Media Integration.}
The \texttt{SocialIngester} module (\texttt{nfl\_analytics/semantic/\allowbreak ingestion/\allowbreak social\_ingester.py}) implements tiered source credibility:

\begin{table}[htbp]
\centering
\caption{Social Media Source Tiers}
\label{tab:social-tiers}
\small
\begin{tabular}{@{} l l S[table-format=1.1] l @{}}
\toprule
\textbf{Tier} & \textbf{Sources} & \textbf{Weight} & \textbf{Examples} \\
\midrule
Tier 1 & Official accounts & 1.0 & @NFL, @Chiefs, @Ravens \\
Tier 2 & Verified insiders & 0.8 & @AdamSchefter, @RapSheet \\
Tier 3 & Beat reporters & 0.5 & Team-specific journalists \\
\bottomrule
\end{tabular}
\end{table}

\paragraph{Historical Training Data.}
The \texttt{AnnotationPipeline} module (\texttt{nfl\_analytics/semantic/\allowbreak training/\allowbreak annotation\_pipeline.py}) generates training data from 84,000+ historical injury records (2009--2025). Key statistics:

\begin{itemize}[itemsep=2pt,parsep=0pt]
    \item \textbf{NER training:} 5,000 samples (4,500 train / 500 validation)
    \item \textbf{Sentiment training:} 2,000 samples from game outcomes (2022--2025)
    \item \textbf{Entity labels:} PLAYER, TEAM, INJURY, STATUS, POSITION
\end{itemize}

\subsection{Layer 2: LLM Inference Architecture}

\paragraph{Design Philosophy: Local-First with Fallback.}
Given a \$50--100/month API budget constraint and available local GPU resources (RTX 4090 on Windows, M4 on Mac), the inference architecture prioritizes self-hosted models with API fallback for reliability.

\paragraph{Fallback Chain.}
The \texttt{SemanticInferenceClient} (\texttt{nfl\_analytics/semantic/\allowbreak models/\allowbreak inference.py}) implements automatic failover:

\begin{enumerate}[itemsep=2pt,parsep=0pt]
    \item \textbf{MLX (Apple Silicon):} Qwen2.5-7B-Instruct-4bit, fastest on M4
    \item \textbf{Ollama (qwen3:8b):} Chat endpoint with 240-second timeout
    \item \textbf{Gemini Flash:} Rate-limited API fallback
    \item \textbf{Claude API:} Final fallback for complex reasoning
\end{enumerate}

\paragraph{Critical Implementation Detail.}
During development (December 2025), we discovered that qwen3:8b's \texttt{generate} endpoint returns empty responses for entity extraction prompts. The solution: use the \texttt{chat} endpoint with structured messages. This endpoint switch resolved 100\% of empty response issues.

\paragraph{Timeout Configuration.}
Complex NER and sentiment queries on qwen3:8b require up to 240 seconds (4 minutes) for thorough analysis. The original 60-second default caused frequent timeouts; increasing to 240 seconds eliminated timeout failures while maintaining acceptable latency for batch processing.

\begin{table}[htbp]
\centering
\caption{LLM Inference Configuration}
\label{tab:llm-config}
\small
\begin{tabular}{@{} l l l S[table-format=3.0] @{}}
\toprule
\textbf{Model} & \textbf{Platform} & \textbf{Endpoint} & \textbf{Timeout (s)} \\
\midrule
qwen3:8b & Ollama & /api/chat & 240 \\
Qwen2.5-7B-Instruct-4bit & MLX & Local & 120 \\
gemini-2.5-flash & API & REST & 60 \\
claude-3-sonnet & API & REST & 60 \\
\bottomrule
\end{tabular}
\end{table}

\subsection{Layer 3: Semantic Feature Engineering}

The \texttt{SemanticFeatureExtractor} (\texttt{nfl\_analytics/semantic/\allowbreak features/\allowbreak semantic\_features.py}) generates 26 features organized into per-team (11 each) and matchup-level (4) categories.

\paragraph{Per-Team Features (11 per team, 22 total).}
Each team receives:

\begin{itemize}[itemsep=2pt,parsep=0pt]
    \item \texttt{sentiment\_24h}: News/social sentiment over past 24 hours ($-1$ to $+1$)
    \item \texttt{sentiment\_7d\_trend}: Week-over-week sentiment change
    \item \texttt{injury\_severity\_llm}: LLM-assessed injury impact (1--5 scale)
    \item \texttt{key\_player\_risk}: Probability of star player absence
    \item \texttt{media\_buzz\_volume}: Normalized media attention (0--1)
    \item \texttt{coaching\_change}: Boolean for recent staff changes
    \item \texttt{roster\_churn}: Player turnover indicator
    \item \texttt{public\_overreaction}: Market sentiment extremity
    \item \texttt{contrarian\_signal}: Inverse public sentiment
    \item Additional derived features (momentum, stability)
\end{itemize}

\paragraph{Matchup Features (4).}

\begin{itemize}[itemsep=2pt,parsep=0pt]
    \item \texttt{sentiment\_differential}: Home minus away sentiment
    \item \texttt{injury\_advantage}: Relative injury situation
    \item \texttt{media\_attention\_ratio}: Coverage imbalance
    \item \texttt{combined\_semantic\_edge}: Composite signal (weighted combination)
\end{itemize}

\paragraph{Integration with GameFeatureExtractor.}
Semantic features are automatically extracted and prefixed with \texttt{sem\_} when integrated into the main feature pipeline:

\begin{verbatim}
features = extractor.extract_game_features(game_id)
# Semantic features: features['sem_home_sentiment_24h'],
#                    features['sem_injury_advantage'], etc.
\end{verbatim}

This prefix convention enables easy identification and ablation of semantic features during model training.

\subsection{Layer 4: Decision Support Interface}

\paragraph{ValueBetDetector.}
The \texttt{ValueBetDetector} (\filepath{nfl\_analytics/production/value\_bet\_detector.py}) identifies value bets where model probability exceeds market-implied probability by a configurable threshold (default: 3\%).

Key capabilities:

\begin{itemize}[itemsep=2pt,parsep=0pt]
    \item \textbf{Edge calculation:} $\text{Edge} = P_{\text{model}} - P_{\text{market}}$
    \item \textbf{Expected value:} EV per \$100 bet
    \item \textbf{Kelly fraction:} Optimal bet sizing (capped at 25\%)
    \item \textbf{Confidence assessment:} HIGH/MEDIUM/LOW based on edge magnitude and probability calibration
\end{itemize}

\paragraph{LLM-Generated Explanations.}
For each value bet, the system generates a 2--3 sentence justification:

\begin{quote}
\textit{``Model identifies 5.2\% edge on Ravens -3.5. Key factor: Strong offensive matchup advantage. Market may be undervaluing Baltimore's rushing attack against Dallas' weakened defensive line.''}
\end{quote}

These explanations are generated via the semantic inference client, providing human-readable context for the human-in-the-loop decision process.

\paragraph{Human-in-the-Loop Design.}
Per user configuration, all bets require human approval. The decision support interface presents:

\begin{enumerate}[itemsep=2pt,parsep=0pt]
    \item Quantitative metrics (edge, EV, Kelly sizing)
    \item Semantic insights (sentiment, injury impact, media signals)
    \item LLM-generated explanation
    \item Risk factors (key player uncertainty, model variance)
\end{enumerate}

This design philosophy---\textit{augment rather than replace human judgment}---aligns with responsible AI deployment in high-stakes financial contexts.

\subsection{Training and Fine-Tuning Pipeline}

\paragraph{Annotation Pipeline.}
The \texttt{AnnotationPipeline} converts database records to training format:

\begin{enumerate}[itemsep=2pt,parsep=0pt]
    \item \textbf{NER samples:} Injury reports $\to$ entity-labeled training examples
    \item \textbf{Sentiment samples:} Game outcomes $\to$ sentiment-labeled examples
    \item \textbf{Export formats:} Ollama JSONL + Modelfile, MLX JSONL
\end{enumerate}

\paragraph{Fine-Tuning Utilities.}
The \texttt{fine\_tune.py} module supports two approaches:

\begin{enumerate}[itemsep=2pt,parsep=0pt]
    \item \textbf{Ollama custom model:} Creates specialized model via Modelfile (system prompt customization, not true fine-tuning)
    \item \textbf{MLX LoRA:} True fine-tuning on Apple Silicon using LoRA adapters (rank=8, 16 layers)
\end{enumerate}

\paragraph{Validation Framework.}
The \texttt{validate\_model} function tests NER performance on held-out cases:

\begin{verbatim}
test_cases = [
    {"text": "Patrick Mahomes questionable with ankle injury",
     "expected_entities": ["Patrick Mahomes", "ankle", "questionable"]},
    {"text": "Chiefs list Travis Kelce as doubtful",
     "expected_entities": ["Chiefs", "Travis Kelce", "doubtful"]},
]
results = validate_model("nfl-ner", test_cases)
# Returns: accuracy, per-case results, missing entities
\end{verbatim}

\subsection{Performance Considerations}

\paragraph{Latency Budget.}
The existing quantitative pipeline maintains $<500$ms end-to-end latency. Semantic features are computed asynchronously:

\begin{itemize}[itemsep=2pt,parsep=0pt]
    \item \textbf{Pre-computation:} Sentiment/injury scores computed every 15 minutes
    \item \textbf{Caching:} Redis cache with content-hash keys
    \item \textbf{Async processing:} Non-blocking LLM calls during feature extraction
\end{itemize}

\paragraph{Cost Management.}
With \$50--100/month API budget:

\begin{itemize}[itemsep=2pt,parsep=0pt]
    \item Local inference (Ollama/MLX): \$0 per query
    \item Gemini Flash fallback: $\sim$\$50/month allocation
    \item Claude API: $\sim$\$25/month for complex reasoning validation
\end{itemize}

\subsection{Empirical Validation: MLX vs Ollama Benchmark}
\label{sec:mlx-ollama-benchmark}

A comprehensive benchmark (December 2025) compared inference performance on Apple M4 Silicon across three task types: injury assessment, sentiment analysis, and named entity recognition. Results strongly favor MLX for production deployment.

\begin{table}[htbp]
\centering
\caption{MLX vs Ollama Inference Benchmark on Mac M4}
\label{tab:mlx-benchmark}
\small
\begin{tabular}{@{} l S[table-format=5.0] S[table-format=4.0] S[table-format=2.1] @{}}
\toprule
\textbf{Metric} & {\textbf{Ollama (ms)}} & {\textbf{MLX (ms)}} & {\textbf{Speedup}} \\
\midrule
Median (p50) & 28627 & 2379 & 12.0 \\
Mean & 31060 & 3323 & 9.3 \\
p95 & 62656 & 7996 & 7.8 \\
Min & 14197 & 536 & 26.5 \\
Max & 62656 & 7996 & 7.8 \\
\bottomrule
\end{tabular}
\end{table}

\paragraph{Task-Specific Performance.}
Sentiment analysis showed the most dramatic improvement, with MLX achieving 23$\times$ faster median latency (723ms vs 16,783ms). NER tasks showed 7.6$\times$ improvement (3,926ms vs 29,719ms). Despite latency differences, throughput remained similar ($\sim$17 tokens/second) due to MLX's batch optimization versus Ollama's streaming approach.

\paragraph{Resource Efficiency.}
MLX consumed approximately 40\% less power (8--12W vs 15--20W) and 33\% less memory (4GB vs 6GB) while achieving superior latency. This efficiency stems from Metal GPU optimization specifically targeting Apple Silicon's unified memory architecture.

\paragraph{Production Recommendation.}
Based on these results, MLX (Qwen2.5-7B-Instruct-4bit) is designated the primary inference backend for Mac deployments, with Ollama retained as cross-platform fallback for CI/CD environments running on Linux.

\subsection{NER Model Evaluation}
\label{sec:ner-evaluation}

The fine-tuned NER model (\texttt{nfl-ner:latest}) was evaluated on a held-out test set of 20 injury report sentences. Overall weighted accuracy reached 76\%, with significant variation across entity types.

\begin{table}[htbp]
\centering
\caption{NER Model Performance by Entity Type}
\label{tab:ner-accuracy}
\small
\begin{tabular}{@{} l l S[table-format=2.0] @{}}
\toprule
\textbf{Entity Type} & \textbf{Performance} & {\textbf{Accuracy (\%)}} \\
\midrule
Player Name & Strong & 95 \\
Team & Strong & 95 \\
Position & Strong & 90 \\
Injury Type & Moderate & 80 \\
\textbf{Status (out/questionable)} & \textbf{Weak} & \textbf{40} \\
\bottomrule
\end{tabular}
\end{table}

\paragraph{Status Extraction as Bottleneck.}
The model reliably extracts player names, teams, and injury types, but struggles with injury status classification (out, questionable, probable, doubtful). In 4 of 5 holdout cases, the model extracted entities like ``ankle injury'' but missed the critical ``questionable'' status. This limitation impacts the accuracy of the \texttt{injury\_severity\_llm} feature.

\paragraph{Implications.}
Status extraction accuracy is a priority for Phase 2 fine-tuning. Current mitigation: weight injury-related semantic features conservatively in ensemble integration until status extraction improves.

\subsection{Section Summary}

The Semantic Stack represents a systematic approach to integrating LLM capabilities into production sports analytics. Key innovations:

\begin{enumerate}[itemsep=2pt,parsep=0pt]
    \item \textbf{Local-first architecture:} Self-hosted models (Ollama, MLX) minimize costs while API fallback ensures reliability
    \item \textbf{38 semantic features:} Sentiment, injury severity/decay, media signals, and derived dynamics (velocity, acceleration) integrated into existing pipeline
    \item \textbf{Human-in-the-loop:} LLM explanations support, rather than replace, human decision-making
    \item \textbf{Training pipeline:} 7,000+ annotated samples from historical NFL data
\end{enumerate}

\textbf{Critical caveat:} A/B testing (Section~\ref{sec:semantic-ab-test}) showed semantic features \emph{underperformed} the baseline on 2025 holdout data. This negative result illustrates the importance of rigorous validation before production deployment. The +1.5--2.5\% win rate uplift target remains aspirational; current evidence does not support enabling semantic features for live betting.

\subsection{A/B Testing: Semantic Features on 2025 Holdout}
\label{sec:semantic-ab-test}

A rigorous A/B test compared spread prediction accuracy with and without semantic features using the 2025 holdout dataset (195 games, Weeks 1--14). Results were \textbf{inconclusive}.

\begin{table}[htbp]
\centering
\caption{Semantic Features A/B Test Results (2025 Holdout)}
\label{tab:semantic-ab-test}
\small
\begin{tabular}{@{} l S[table-format=2.1] S[table-format=2.1] @{}}
\toprule
\textbf{Model} & {\textbf{Spread Accuracy (\%)}} & {\textbf{Spread MAE}} \\
\midrule
Control (no semantic, v1) & 58.5 & 10.2 \\
Treatment (with semantic, v2) & 51.8 & 10.8 \\
\midrule
\textbf{Difference} & {-6.7} & {+0.6} \\
\textbf{McNemar p-value} & \multicolumn{2}{c}{0.21 (not significant)} \\
\bottomrule
\end{tabular}
\end{table}

\paragraph{Interpretation.}
The semantic-enhanced model (v2) underperformed the baseline on this holdout set, though the difference was not statistically significant ($p = 0.21$, McNemar test). Several explanations warrant investigation:

\begin{enumerate}[itemsep=2pt,parsep=0pt]
    \item \textbf{Feature engineering quality:} NER status extraction weakness (Section~\ref{sec:ner-evaluation}) may inject noise into injury severity features
    \item \textbf{Distribution shift:} 2025 data patterns may differ from 2022--2024 training distribution
    \item \textbf{Signal dilution:} 26 additional features may overwhelm rather than complement quantitative signals
    \item \textbf{Sample size:} 195 games may be insufficient to detect a 1--2\% true effect size
\end{enumerate}

\paragraph{Lessons Learned.}
This result validates the importance of rigorous A/B testing before production deployment. Had semantic features been enabled based solely on theoretical appeal or training set performance, they would have degraded 2025 predictions. The negative result is documented here in the spirit of ``honest failure accounting'' (Section~\ref{sec:honest-failures}), which argues that null or negative results are essential for building credible research.

\paragraph{Future Work.}
Phase 2 priorities, informed by empirical validation:
\begin{itemize}[itemsep=2pt,parsep=0pt]
    \item \textbf{RAG integration:} Historical context retrieval for situational awareness
    \item \textbf{NER status improvement:} Additional training data focused on status extraction (out/questionable/probable)
    \item \textbf{Feature ablation:} Identify which semantic features (if any) provide positive signal
    \item \textbf{Extended A/B testing:} Continue monitoring through 2025 Week 18 for larger sample size
\end{itemize}
